%% -----------------------------------------------------
%% Sample free form document written in NSERC format
%%
%% Declare the document as "nserc-appendix"; use either 11 or 12pt
%% fonts, and specify whether the document is in French or English
%% -----------------------------------------------------
\documentclass[11pt,english,calcwidth,compact]{nserc-appendix}

%% --------------------------
%% Inclusion of a basic configuration file. Go see this file: there are
%% parameters (such as your name, etc.) to fill in. This avoids repeating
%% the same info in the case you produce multiple documents for the same
%% application.
%% --------------------------
\input{application-config.tex}
\usepackage[compact]{titlesec}
\usepackage[protrusion=true,expansion=true]{microtype}
%\usepackage{overcite}
%% ----------
%% Set the title of the document in the PDF's metadata
%% ----------
\hypersetup{%
  pdftitle = {AI-Driven Quantum Chemistry for Molecular Materials}
}

\begin{document}
\title{AI-Driven Quantum Chemistry for Materials Chemistry}

%% ----------
%% Set start page number of the document. Change this depending on
%% the part you are writing, and its position relative to the rest of
%% the application.
%% ----------
\setcounter{page}{1}
\maketitle

This research program aims to overcome the long-standing challenge of accurately modelling complex molecular and solid-state systems by developing an integrated suite of computational tools that combine dispersion‐corrected density-functional theory (DFT) with cutting-edge artificial intelligence (AI) techniques. In doing so, it will bridge theoretical insights with experimental validation, ultimately enabling high-throughput predictions of materials properties. The proposed methods will be particularly relevant for designing semiconductors, two-dimensional (2D) materials, and defect-rich solids, which play pivotal roles in fields as diverse as nuclear chemistry, minerals processing, and optoelectronics. By systematically improving both the accuracy and efficiency of DFT calculations, as well as integrating experimental observables into the computational workflow, this research lays the foundation for accelerated discovery of next-generation functional materials.

\section{Background and Objectives of Research Program}

Over the past two decades, dispersion-corrected DFT\cite{jcp-perspective,becke2014JCP,burke2012JCP} has emerged as a central tool for computational materials chemistry, enabling predictions of crystal structures \cite{moellmann2010importance, mayo2024assessment,BT7Rank}, electronic properties \cite{grimme2016dispersion}, and intermolecular interactions \cite{xdmaims}, with moderate computational cost. However, conventional DFT methods still struggle in several key areas. For instance, they often require extensive parameter sweeps \cite{gould2022poisoning,PhysRevX.15.011013,langreth1975exchange}, struggle with delocalization error \cite{wires}, and can underperform for materials containing localized electronic states \cite{yao2015efficient}, such as transition metal d-orbitals or defects, like vacancies and dopants \cite{sabino2022intrinsic}. My doctoral research contributed to addressing these issues in molecular crystals, successfully applying dispersion-corrected DFT to predict polymorphism via accurate modelling of intermolecular interactions in organic solids \cite{xdmaims,aimscsp}. Building on that foundation, my postdoctoral work introduced adaptive DFT concepts \cite{apbe0}, where machine learning (ML) algorithms dynamically optimize simulation parameters (such as the degree of exact-exchange mixing in hybrid DFT), thereby reducing computational overhead and achieving higher fidelity in gas-phase calculations for organic molecules.

A complementary development is a machine learning framework called “amons” \cite{amons}, originally devised for molecular systems, which decomposes large structures into local fragments (amons) that capture critical bonding environments. I am now extending this approach to solid-state amons, aiming to describe extended networks and realistic defect concentrations in materials such as oxides and layered semiconductors. This adaptation allows for modelling far larger systems than would be feasible with standard supercell DFT. By integrating ML models trained on carefully selected amons, we gain predictive power across broad chemical spaces, from electronically active dopants in wide-gap oxides to surface defects in mining- and catalysis-relevant minerals.

The principal objective of my proposed research is to create an end-to-end computational framework that unites adaptive DFT strategies with solid-state amons and experimental data. In the short term, my group will develop \textbf{double hybrid and adaptive machine learning for DFT}, where generation and leverage of high-level reference data will be used to refine dispersion-corrected DFT on-the-fly, reducing reliance on costly parameter sweeps while improving accuracy for challenging materials. In parallel, we will develop \textbf{solid-state amons for defect-rich materials}, where the amons concept is extended to complex solids, enabling accurate yet computationally tractable modelling of low-concentration defects, doped semiconductors, molecular materials, and 2D materials. This approach will circumvent the size limitations of conventional plane-wave supercell calculations, making realistic defect studies (e.g., 0.1\% doping) feasible. The third objective is \textbf{integration with experimental observables}, where we will develop protocols to incorporate real-world measurements (such as lattice constants, defect formation energies, and chemical shifts) into the model training process, creating a feedback loop that tightly couples theory and experiment.

My long-term vision is to  catalyse progress in materials chemistry on multiple fronts by developing a comprehensive suite of tools that seamlessly integrate AI with dispersion-corrected solid-state DFT in a streamlined workflow. This approach aims to achieve high predictive accuracy at efficient computational costs, thereby accelerating the discovery and design of advanced materials. 
The tools we develop can dramatically reduce the time and expense associated with high-level, solid-state simulations, opening doors to more precise modelling of defect-related phenomena in semiconductors, catalysts, and nuclear materials. These capabilities, in turn, have immediate implications for energy technologies, environmental sustainability, and the fundamental understanding of material properties. Ultimately, the convergence of computational efficiency, experimental feedback, and advanced AI promises to transform how researchers discover and optimize new materials, paving the way for breakthroughs in electronics, photonics, energy conversion, and beyond.

Although my doctoral research established a strong track record in molecular crystals, this proposal expands into new material classes (semiconductors, 2D systems, and defect-laden solids) and introduces AI-driven techniques to meet the growing demand for predictive modelling in energy, mining, and advanced manufacturing contexts.
I bring a unique combination of expertise in dispersion-corrected DFT for molecular solids, algorithmic innovations in adaptive DFT from my postdoctoral work, and experience incorporating AI methods into traditional quantum chemistry. These skills directly align with the challenges identified above—namely, bridging quantum accuracy with large-scale materials design. Further, my established collaborations with experimental scientists will facilitate the acquisition of high-fidelity observables for training and validating our computational models.


%------------------------OLD STUFF -----------

%\section*{old stuff}
%This research program aims to overcome the long-standing challenge of accurately modelling non-covalent interactions in materials chemistry by developing an integrated suite of computational tools that combine dispersion‐corrected solid‐state DFT with cutting-edge AI and machine learning techniques. Building on a strong foundation in crystal structure prediction for molecular solids, the program will extend these methodologies to encompass semiconductors, two-dimensional systems, and defect-rich materials. In the long term, the goal is to create a high-throughput, predictive framework that not only reduces the computational cost of high-fidelity simulations, but also bridges theoretical insights with experimental observables through adaptive ML algorithms and a self-driving laboratory platform. In the near term, the work will focus on optimizing DFT parameters via dynamic ML models, expanding the concept of solid-state atoms in molecules to deconstruct complex materials into manageable subunits, and automating synthesis and characterization workflows to enable rapid feedback and refinement. This holistic approach promises to accelerate the discovery and design of next-generation functional materials for advanced applications in electronics, nuclear energy, energy conversion, and beyond.
%
%
%\section{Background and Objectives of Research Program}
%
%The accurate treatment of non-covalent interactions remains a central challenge in modern computational materials chemistry. Over the past decades, dispersion‐corrected density functional theory (DFT) has emerged as an indispensable tool for the prediction of intermolecular interactions and enabling reliable crystal structure prediction (CSP) of molecular solids. My research career has been anchored in the development and application of dispersion‐corrected solid‐state DFT techniques, with a particular focus on CSP for molecular solids. In this context, the ability to predict the relative stabilities of various polymorphs from the chemical structure of a monomer has profound implications for materials design and synthesis. However, while notable successes have been achieved in molecular CSP, a vast territory of unexplored systems remains. In particular, semiconductors, two-dimensional (2D) materials, and materials containing complex defect structures --- all of which have wide applications in the fields of minerals and mining chemistry, nuclear chemistry, and surface science --- stand to benefit immensely from enhanced computational methodologies.
%
%My \uline{long-term} vision is to advance materials chemistry by developing a comprehensive suite of tools that seamlessly integrate artificial intelligence (AI) with dispersion-corrected solid-state DFT, and experiment. This approach aims to achieve high predictive accuracy at efficient computational costs, thereby accelerating the discovery and design of advanced materials. Building on my established expertise in CSP of molecular solids, I plan to extend these methods to explore a broader range of materials. In doing so, the research program will bridge theoretical advances with practical experimental results, establishing a workflow for the predictive design of solid-form functional materials.
%
%
%In the \uline{short and near term}, I will focus on three interconnected objectives that collectively underpin this framework. First, I will develop adaptive machine learning (ML) algorithms to optimize computational parameters for solid form materials by integrating high-level reference data with experimental observations, such as single crystal and powder X-ray diffraction data, RPA-DFT, and NMR observables. This integration is crucial for balancing computational cost with accuracy, enabling rapid yet reliable predictions of intermolecular interactions and material properties. Second, I will expand the concept of ''solid-state Atoms in molecules (AMONs)''---previously applied to molecular systems---into the solid-state domain. By systematically identifying and utilizing representative substructures that capture the essential chemical features of complex materials, I aim to simplify the exploration of materials and the behaviour of defects. Finally, I intend to apply these computational advances within a self-driving laboratory framework. By establishing a closed-loop system that integrates real-time computational predictions, experimental synthesis, and Bayesian optimization to identify gaps in the current data, this approach will accelerate the discovery and optimization of novel solid-form materials.
%
%While my PhD work provided a strong foundation in molecular crystal structure prediction, my research trajectory is distinguished by its expansion into new material classes and AI driven methodological innovations. By bridging traditional dispersion-corrected DFT with cutting-edge AI techniques, I intend to develop versatile tools that address current challenges in materials design. In particular, my focus on semiconducting, 2D, and defect-rich materials represents a significant evolution from my earlier work, opening avenues for transformative advances in electronics, optoelectronics, and energy conversion \cite{DBLPjournals/jsyml/Turing48,DBLPconf/afips/SolomonP76}.
%
%%\section{Long-Term Objectives}
%
%At the core of my long-term research objectives is the ambition to further refine the computational treatment of intermolecular chemistry by establishing a framework that seamlessly integrates dispersion-corrected solid state DFT with AI-driven methodologies. The primary goal is to develop a suite of computational tools capable of accurately predicting crystal structures, electronic properties, and defect behaviours across a wide array of materials systems, ranging from traditional molecular solids to more complex condensed-phase materials, and providing these tools within a readily accessible package to the larger chemistry community. 
%
%To accomplish this, I intend to leverage the inherent strengths of dispersion-corrected DFT, its capacity to capture subtle electronic interactions, while mitigating its computational expense and known errors through the integration of adaptive ML algorithms. These algorithms will learn from both high-level reference data and experimental inputs, thereby reducing the need for exhaustive parameter sweeps and enabling efficient, and high-throughput screening of candidate materials. Additionally, by expanding the conceptual framework of “solid-state amons” to include semiconducting and 2D materials, I will develop a modular approach that deconstructs complex systems into representative building blocks, allowing for the investigation of large physically meaningful catalytic surfaces, and energetic materials. This strategy will facilitate targeted investigations into the interplay of dispersion forces, electronic structure variations, and defect chemistry in determining material performance, ultimately paving the way for the design of next-generation materials with tailored properties.
%
%%\section{Concluding Remarks}
%
%In summary, my proposed research program aims to advance the state-of-the-art in computational materials chemistry by developing a comprehensive suite of tools that integrate dispersion-corrected solid‐state DFT, adaptive machine learning, and automated experimental platforms. Through detailed algorithm development, the expansion of solid-state amons, and the realization of a self-driving laboratory, this program seeks to address longstanding challenges in materials design.
%
%By differentiating my work from previous research and extending the focus to include semiconducting, 2D, and defect-rich materials, the program not only builds on my established expertise but also paves the way for novel applications with significant technological impact. The successful execution of this research will provide the computational chemistry community with robust, efficient, and versatile tools for exploring complex materials, thereby facilitating high-throughput screening and rational design. This, in turn, will enable the discovery of new materials with tailored properties for applications in electronics, photonics, energy conversion, and beyond.
%
%Ultimately, the integration of AI-driven methodologies with state-of-the-art computational chemistry stands to transform our approach to materials research, leading to breakthroughs with far-reaching implications across science and industry.
%
\section{Proposed Research}

\subsubsection*{Theme I: Adaptive Density Functional Theory for Solids}

This theme is dedicated to enhancing the efficiency and accuracy of dispersion-corrected DFT calculations through the integration of adaptive machine-learning (ML) techniques and their incorporation into current electronic structure workflows \cite{apbe0}. Adaptive DFT for solids represents a shift in computational materials chemistry, offering routes to address critical challenges such as delocalization and basis-set superposition errors, while preserving or lowering computational costs associated with traditional methods. In adaptive DFT, ML techniques are integrated to dynamically optimize the computational parameters during simulations, allowing the method to tailor itself to the specific chemical environment of a system. The dynamic optimization enables more accurate predictions of both bonded and non-bonded interactions, which is particularly crucial for complex solid-state systems where a delicate balance of forces determines the material properties. By leveraging high-level reference data in the form of double-hybrid functionals applied to solids,  adaptive DFT will be able to refine its parameters on-the-fly, ultimately bridging the gap between theoretical accuracy and practical computational expense.  

\subsubsection*{Project I(a): Implement Improved Double Hybrids for Solids}  

This project aims to implement and optimize double-hybrid (DH) functionals \cite{dh23} specifically tailored for solid-state systems, thereby significantly enhancing the accuracy of DFT calculations for complex materials. Double-hybrid functionals combine conventional DFT with a perturbative correction from Møller–Plesset perturbation theory (MP2), offering a more rigorous treatment of electron correlation effects \cite{dh1,dh2}. However, their application to solids has been limited due to the prohibitive computational cost of traditional plane-wave–based methods. By adopting an efficient numerical atom-centred orbital framework provided by the FHI-aims code \cite{fhiaims,fhiaimshybrid}, our implementation will reduce the computational cost associated with MP2-like calculations and enable the optimized selection of smaller basis sets. This approach not only improves the intrinsic accuracy of the method but also fills a critical gap in the field, as similar implementations are not currently being pursued elsewhere.

The optimized DH framework will serve as a robust tool for generating high-quality computational data for solids --- data that have traditionally been restricted to hybrid and generalized gradient approximation levels. The double-hybrid approach will be essential both for predictive modeling of solids and for constructing training datasets for future methodological developments. The enhanced performance will facilitate accurate simulations of large and complex solid-state systems that are currently beyond the reach of conventional methods. Examples include defect-laden materials and extended catalytic surfaces that encompass more than just the active site. This project, therefore, represents a significant methodological advance in computational materials chemistry by enabling more realistic and cost-effective modelling of solids. This subproject will be undertaken by one PhD student and will provide training within a well-established code base that is widely used, both domestically and internationally. 

\subsubsection*{Project I(b): Implementation of Adaptive DFT for solids}  

In this project, I will develop adaptive machine learning algorithms that optimize the computational parameters of dispersion-corrected solid-state DFT calculations. By systematically integrating high-level reference data, the ML models will be capable of dynamically adjusting key simulation parameters depending on the chemical system under study. This approach was developed during my post-doctoral work, has shown great promise for gas-phase calculations, and will ensure that calculations maintain high accuracy while operating at reduced computational cost relative to the reference levels. A particular focus will be placed on resolving the fine balance among competing intermolecular interactions, such as hydrogen bonding, $\pi$-stacking, van der Waals forces, and delocalization error, which is critical for reliable solid-state calculations. Initially, the adaptive algorithms will be validated against well-defined benchmark systems, such as a set of 55 diverse benchmarks for main-group thermochemistry, reaction barriers, and non-covalent interactions (GMTKN55) \cite{gmtkn55}, a set of 23 molecular crystal lattice energies (X23) \cite{x23,x23b}, and a set of 110 million DFT calculations 
%focused on structural and compositional diversity 
of inorganic bulk materials (OMat24) \cite{omat24}, before being applied to more complex cases involving semiconducting and 2D materials. 
Implementing adaptive DFT in solids relies on an ML pipeline capable of generalizing from limited training data. To this end, we adopt the many-body distribution functional (MBDF) representation, which encodes local atomic environments using mathematical functionals of two- and three-body distributions \cite{mbdf,cmbdf}. Given the sparse data typically available for complex solids, we will employ kernel ridge regression (KRR) \cite{vapnik1999nature,rasmussen2006gaussian, GPR_deringer} to predict optimal DFT exchange-correlation parameters (e.g., the exact-exchange mixing fraction \cite{apbe0}) with minimal overhead. 
%As the database of solid-state reference calculations grows, we will explore neural network architectures—such as deep density functional models—to capture more intricate quantum relationships. 
This research will thus form a robust platform for future methodological advances and will serve as a compelling training ground for 1–2 graduate students, who will gain expertise in computational chemistry, ML approaches, and solid-state modelling.
%Implementing adaptive DFT in solids relies on efficient machine learning models that can generalize from limited training data. We adopt the Many-Body Distribution Functional (MBDF) representation introduced by Danish Khan et al. as a compact way to describe atomic environments. Briefly, MBDF uses mathematical functionals of two-body and three-body atomic distributions (and their derivatives) to characterize each atom’s local chemical environment. %This representation, which can be made convolutional for periodic systems (cMBDF), provides a fixed-length feature vector per structure capturing both short-range and medium-range geometric information. 
%Given the relatively sparse data available (especially for complex solids), we use a Kernel Ridge Regression (KRR) model as the machine-learning engine. %KRR is well-suited for low-data regimes due to its strong performance with limited training examples. 
%In our adaptive functional, KRR learns to predict the optimal DFT exchange-correlation parameter (e.g. the exact exchange mixing fraction) from the MBDF features of a given system. This data-efficient, nonparametric approach yields reliable predictions of the functional tuning “on the fly” with minimal overhead. 
%As our database of reference calculations grows (e.g. incorporating more solid-state training points), we can consider transitioning to more flexible neural network models. Neural networks (such as deep density functional models) have shown promise in capturing complex quantum relationships once sufficient data is available. The insights obtained from this project will provide a firm foundation for future methodological innovations. This project would be conducted by 1-2 first year graduate students, providing them with foundational training in computational chemistry, machine learning techniques, and solid-state modelling.  

These projects ultimately aim to integrate state-of-the-art adaptive hybrid density functionals (aDFT) in addition to novel DH functionals into established electronic structure packages, specifically FHI-aims. The incorporation of functionals designed to address delocalization errors, and enhance the treatment of exchange-correlation effects, within  FHI-aims will improve the reliability and efficiency of DFT calculations key to future AI-based training. 
By embedding these functionals into an efficient computational workflow, the project will facilitate accurate predictions for a broad range of systems, particularly in the solid-state regime.  
Benchmarking against standard datasets (GMTKN55 \cite{gmtkn55}, X23 \cite{x23,x23b}, VQM24 \cite{vqm24}, and QM9 \cite{qm9}) will ensure that the newly incorporated functionals perform robustly, thereby enabling their application in complex materials such as defect-rich materials and nuclear materials. 
%The implementations of these subprojects will provide training within a well established, and widely used code base domestically and internationally, for the graduate students involved.  

\subsubsection*{Theme II: Application of ML Methods to Solid-Form Defect Materials}

In this theme, we will extend the atoms-in-molecules (amons) framework \cite{amons}, originally developed for molecular systems, to solid-state materials chemistry. This approach relies on the intelligent selection of progressively larger supercells, where each configuration serves as a representative fragment that captures essential local defect chemistry and electronic structure. We will fully integrate the amons-based approach into the FHI-aims package \cite{fhiaims}, creating a cohesive workflow for high-throughput screening of defect-rich systems. Given a bulk crystal structure and defect specification (e.g., a dopant in a layered TMDC), the software will automatically generate  at set of local fragments (amons) and perform rigorous quantum-mechanical calculations on those fragments. By focusing on smaller, parallelized computations, we can afford advanced functionals, such as hybrids or double-hybrids, for near-benchmark accuracy. Each amon fragment then contributes to a shared library of defect configurations, which an ML model uses to predict properties (dopant positions, concentrations, or strain effects) in larger systems without recalculating from scratch. This modular strategy enables detailed insights into doping asymmetries, exciton recombination, and oxygen diffusion phenomena, ultimately accelerating both fundamental research and practical materials engineering.

Many metal oxides (e.g., ZnO, SnO$_2$, NiO, Ga$_2$O$_3$, Gd$_2$O$_3$)  serve in applications ranging from transparent electronics to radiation-resistant coatings \cite{transparentcoat,nucshield}, but pose notable modelling challenges due to strong d-electron localization and complex defect chemistry. Semi-local DFT often underestimates band gaps and misplaces localized defect states, yielding qualitatively incorrect results at realistic doping levels \cite{localized2011defect}. Layered transition-metal dichalcogenides (TMDCs), like MoS$_2$ or WS$_2$, similarly involve correlated d-states and can form intricate defects such as vacancies, grain boundaries, or charge-density waves \cite{kim2022experimental}. Modelling such low-concentration defects with conventional supercell DFT can be prohibitively expensive, especially when large simulation cells are needed to avoid artificial interactions between periodic images. By combining the amons framework with efficient ML-driven workflows, we aim to tackle these computational bottlenecks, providing accurate and scalable solutions for defect-rich oxides and 2D materials.
%
%We extend this principle to “solid-state amons” by systematically exploring a series of defected supercells—starting from smaller cells that capture the local environment around a dopant or vacancy—and progressively scaling up. Each supercell effectively serves as a training example for a specific local motif (e.g., a dopant in a particular coordination). By varying defect positions, doping concentrations, or supercell sizes, we generate a library of configurations that collectively resembles the modular “amons” approach in molecular systems. Once this training dataset has been amassed, we can predict properties of much larger or more complex supercells without explicitly recalculating every atomic interaction. In practice, this technique enables efficient modelling of dilute dopants, extended defect clusters, and other challenging scenarios where a brute-force first-principles simulation would be prohibitive.
%
%We will fully integrate this amons-based approach into the FHIaims electronic structure package, creating a cohesive workflow for automatic fragment generation, property evaluation, and high-throughput screening of defect-rich materials. Given a bulk crystal structure and defect specification (e.g., a substitutional dopant in a TMDC monolayer), the workflow will automatically extract the relevant local cluster, and perform rigorous quantum calculations on the resulting fragment. By enabling parallelized, smaller-scale computations, we can feasibly employ advanced functionals—such as hybrids or double-hybrids—to achieve near-benchmark accuracy. Each amon fragment contributes to a growing library of defect configurations, which in turn trains machine-learning models capable of predicting properties for new variants—such as different dopant positions, concentrations, or strain conditions—without recalculating the entire system from scratch. This modular and data-driven framework will shed light on critical defect-driven phenomena, including doping asymmetries in wide-gap oxides, exciton recombination in layered TMDCs, and oxygen diffusion pathways in nuclear materials, thereby accelerating both fundamental research and practical materials engineering.
%
%Many metal oxides (e.g., ZnO, SnO₂, NiO, Ga₂O₃) are critical in minerals processing, catalysis, and even nuclear materials. They often serve as wide-band-gap semiconductors in sensors, transparent electronics, and radiation-resistant coatings. These oxides present modelling challenges because of strongly localized d-electrons and pronounced defect chemistry. For instance, wide-gap oxides like ZnO and SnO₂ are easily n-type doped while p-type doping is extremely difficult, reflecting a “doping asymmetry” that conventional DFT struggles to capture. Standard semi-local DFT tends to underestimate band gaps and fail in describing localized defect states (e.g., small polarons in oxides), leading to qualitative errors in defect level positions and formation energies. This makes it hard to predict real-world behaviour (such as conductivity changes at low dopant levels) with traditional methods alone. 
%Layered Transition Metal Dichalcogenides (TMDC) materials (e.g., MoS₂, WS₂, SnS₂) are another focus due to their importance in surface science, 2D electronics, and mining chemistry (molybdenite MoS₂ is a common mineral and an industrial catalyst). These compounds have a 2D layered structure, meaning their surfaces and interfaces dominate many properties. Defects in TMDCs (vacancies, substitutions, grain boundaries) can drastically alter optical and electronic behaviour, but accurately modelling them is challenging. Conventional simulations require large supercells to isolate a single defect in a sheet, and even then the long-range effects (e.g. strain fields or interlayer interactions) can be missed. Moreover, TMDCs involve transition-metal d orbitals and possible charge-density wave or excitonic effects, which push the limits of standard DFT. In short, both oxide semiconductors and TMDCs tend to be computationally demanding: they contain complex, correlated electronic states and require accounting for dilute defect populations, conditions under which conventional methods can be unreliable or prohibitively expensive.
%
\subsubsection*{Project II(a): Expansion of Current Methods to Create a Defect-Focused Materials Dataset}  

We extend the amons concept to solids by generating a series of defect-containing supercells that progressively increase in size and complexity, each capturing a specific local environment around a vacancy or dopant. Taken together, these supercells form a robust library of representative amons that can be used to train ML models. This modular dataset allows us to predict properties of much larger or more complex supercells, including those with low-level doping or extended defects, without recalculating every atomic interaction. By systematically varying defect positions, concentrations, and cell dimensions, we effectively encode the key chemical motifs behind defect formation and propagation, all while avoiding the exponential cost of brute-force simulations.

Implementation within the FHI-aims package will support efficient property evaluation and high-throughput screening of defect-rich materials. Smaller-scale computations on these amons can be parallelized, enabling the use of advanced functionals (e.g., hybrids or double-hybrids) for near-benchmark accuracy. Each amon entry feeds into a growing ML model that accurately interpolates to different dopant positions, concentrations, or strain states—crucial for tackling phenomena like doping asymmetries in wide-gap oxides or exciton formation in layered TMDCs. This approach also alleviates the common pitfall of artificially high defect concentrations in conventional supercell DFT, enabling realistic studies of, for example, 0.1\% substitutions in large crystals. Over time, a dedicated graduate student will expand the library with additional materials and defect configurations, culminating in a publicly accessible resource that accelerates both fundamental and applied research on defect-laden solids.
%To overcome these challenges, I propose using amons --- ``atoms-in-molecules'' fragments --- adapted for solids. Originally, amons were developed in molecular quantum machine learning as small, representative molecular fragments that capture the essential local chemistry of larger molecules, substituting an atomic basis with a molecular basis. The idea is that larger structures can be reconstructed from local building blocks: if you include all relevant local environments (bonding motifs, functional groups) as separate fragments, you can predict the properties of a complex molecule with high accuracy by learning from those fragments. In the molecular domain, an amons-based model actively selects minimal fragments that reproduce a molecule’s local atomic environments, allowing quantum-accurate ML predictions with far less data.
%We extend this principle to “solid-state amons” by systematically exploring a series of defected supercells—starting from smaller cells that capture the local environment around a dopant or vacancy—and progressively scaling up. Each supercell effectively serves as a training example for a specific local motif (e.g., a dopant in a particular coordination). By varying defect positions, doping concentrations, or supercell sizes, we generate a robust library of configurations that collectively resembles the modular “amons” approach in molecular systems. Once this training dataset has been amassed, we can predict properties of much larger or more complex supercells without explicitly recalculating every atomic interaction. In practice, this technique enables efficient modeling of dilute dopants, extended defect clusters, and other challenging scenarios where a brute-force first-principles simulation would be prohibitive.

%We will fully integrate this amons-based approach into the FHIaims electronic structure package, creating a cohesive workflow for property evaluation and high-throughput screening of defect-rich materials. Rather than physically cutting out gas-phase fragments, the method systematically constructs smaller supercells that capture the local environment around each defect or dopant. By progressively scaling supercell size and defect configurations, we generate a diverse library of “amons” that collectively represent the underlying chemistry of the larger crystal. These smaller-scale computations can be parallelized, allowing us to employ advanced functionals (e.g., hybrids or double-hybrids) without incurring prohibitive computational costs. Each amon entry contributes to a growing dataset used to train machine-learning models, enabling rapid predictions for new variants—such as different dopant positions, concentrations, or strain states—without recalculating every atomic interaction. This modular and data-driven framework will illuminate critical defect-driven phenomena, including doping asymmetries in wide-gap oxides, exciton recombination in layered TMDCs, and oxygen diffusion pathways in nuclear materials, thereby accelerating both fundamental research and practical materials engineering. This project will be an ongoing project by 1 graduate student, where the library of training data will be increased as subsequent trainees are added to the group. A major advantage of the amons framework is the ability to simulate realistic defect concentrations and dopant loads that are currently out of reach with direct first-principles modeling. In traditional supercell calculations, a single impurity or vacancy in a reasonably sized cell might correspond to an unrealistically high doping concentration (often several atomic percent). For example, modeling an experimentally dilute doping (say 0.1\% substitution) would require thousands of atoms in a cell – an unfeasible supercell size for brute-force DFT. 

\subsubsection*{Project II(b): Using Amons to Study Oxygen defects}  
%\textbf{Incorporate into known workflow}  
 
 Oxygen vacancies in oxides critically influence electronic structure, impacting doping, conductivity, and catalytic properties. For instance, an oxygen vacancy generates Ti$^{3+}$ centers and in-gap states in TiO$_2$ \cite{bharti2016formation}, whereas it induces shallow donor levels near the conduction band edge in ZnO \cite{KIM20121711}. Similarly, in UO$_2$, each vacancy can reduce U$^{4+}$ to U$^{3+}$, altering thermal conductivity and oxidation state \cite{dorado2012first}. Such phenomena are central to semiconductor performance, catalytic reactivity, and defect evolution in nuclear fuels, but accurately modelling them remains challenging.

Conventional plane-wave DFT often underestimates oxide band gaps, struggles with self-interaction errors, and relies on large supercells to simulate dilute defects, driving up computational costs. Methods like DFT+U \cite{wang2006oxidation,aykol2014local,macke2024orbital}, hybrid functionals \cite{gruneis2014ionization, chen2014band}, or GW \cite{chen2017accuracy} can localize defect charges and recover the correct band gap, but each introduces parameter sensitivities or further raises the computational burden. The amons-based approach addresses these limitations by modelling only the local chemical environment around the vacancy.  By building up smaller supercells, we obtain compact systems amenable to higher-level methods, providing high-accuracy data on defect-induced electronic and structural perturbations, which can then be used to train predictive ML models. Machine learning models trained on these subunits can then predict defect formation energies, ionization levels, and polaron states in the extended solid, bypassing the need for massive supercell calculations. This strategy thus balances accuracy and scalability, enabling first-principles insights into oxygen vacancies under realistic doping conditions. In essence, the amons approach uses transferable local building blocks to reconstruct accurate electronic structure for defective oxides, achieving periodic-quality predictions for band gaps and defect states in complex oxides at a fraction of the computational cost.

% By building up smaller supercells, we obtain compact systems amenable to higher-level methods capable of providing accurate data  \textcolor{red}{encoding defect–host interactions --- not sure what you mean here}. 

%Substituted 


 %Oxygen vacancies are fundamental defects in oxides that strongly perturb the electronic structure. Removing an oxygen leaves behind electrons that can shift band edge positions or occupy mid-gap states. For example, in TiO₂ an oxygen vacancy creates Ti³⁺ centers with a defect state about 1 eV below the conduction band, and in ZnO it acts as an n-type dopant by introducing shallow donor levels near the conduction band edge. In a correlated oxide like NiO, oxygen vacancies inject extra electrons into Ni d-states, producing localized mid-gap states that cause sub-bandgap optical absorption (colour centers). Likewise, in UO₂ each neutral oxygen vacancy leaves two electrons that reduce U⁴⁺ to U³⁺, altering the local charge balance and electronic conductivity. Understanding these vacancy-induced effects is crucial across applications: in semiconductors they control doping and charge carriers, in catalysis they often serve as active sites (the ease of forming surface O vacancies correlates with catalytic activity), and in nuclear materials like UO₂ the oxygen vacancy concentration influences oxidation state, thermal conductivity, and defect evolution in fuel. 
 %Simulating oxygen defects with standard plane-wave DFT is notoriously difficult. Semilocal DFT (e.g. PBE) severely underestimates oxide band gaps and tends to delocalize the vacancy’s electrons, often predicting no in-gap state and a Fermi level pinned in the conduction band. This failure is due to self-interaction error causing poor charge localization---for instance, PBE may smear the two vacancy electrons across the lattice instead of forming a localized small polaron on a metal site. Correcting this requires beyond-DFT treatments: DFT+U, hybrid functionals, or GW can recover the band gap and localize defect charge, but each comes with trade-offs. DFT+U depends on the chosen U parameter (different U values alter the degree of localization and defect level positions), while hybrid and GW methods are far more computationally expensive. Moreover, modelling a dilute oxygen vacancy demands a very large supercell (hundreds of atoms) to minimize spurious defect–defect interactions. If the cell is too small, the artificially high defect concentration can distort results --- for example, in TiO₂ the stable polaronic state changes when the supercell size is reduced, due to interaction between periodic vacancies. These size and method requirements make conventional plane-wave DFT calculations for isolated vacancies computationally prohibitive. The amons-based approach offers a compelling solution: instead of explicitly simulating a huge periodic lattice, one uses small fragment models (“amons”) that capture the local chemical environment of the defect. By exploiting the locality (nearsightedness) of electronic interactions, an amons framework can include the defect and its nearby cations, yet still mimic the screening and relaxations as in the bulk solid. Quantum calculations on such subunits --- chosen systematically to represent all relevant defect chemistry --- provide training data that encode defect–host interactions and electronic screening effects in a manageable size. Machine learning models trained on these amons can then predict the defect’s influence on the extended material, such as band edge shifts due to vacancy doping, defect ionization energies/levels, and charge-trapping or polaron formation energies, all without requiring an explicit large supercell. In essence, the amons approach uses transferable local building blocks to reconstruct accurate electronic structure for defective oxides, achieving periodic-quality predictions for band gaps and defect states in complex oxides at a fraction of the computational cost.

\subsubsection*{Theme III: Bridging the divide between experimental observables and computational results}

Integrating experimental observations with computational modelling is crucial to enhance the realism and accuracy of materials simulations. Current ML models for material properties---including adaptive DFT schemes and the amons approach---are predominantly trained on in silico data (e.g. DFT calculations), which limits their direct applicability to real-world conditions. To truly bridge the divide, we will incorporate high-fidelity experimental data at multiple stages of model development. This strategy will calibrate ML predictions to reflect actual material behaviour and not just reproduce computational biases.

\subsubsection*{Project III(a):Experimental Data validation and use.}  
The overarching goal of this theme is to incorporate high-fidelity experimental data into our computational workflows, ensuring that the ML models guiding both adaptive DFT and solid-state amons remain tethered to real-world behaviour. Initially, we will rely on published datasets—such as well-characterized crystal structures (from the ICSD \cite{levin2018nist} or COD \cite{COD,cod2}) and documented nuclear magnetic resonance (NMR) measurements—to create a reference set of material observables. These data will serve as benchmarks for calibrating the accuracy of our computational methods, revealing where predictions from both standard DFT and ML-based approaches deviate significantly from experiment. By incorporating even a limited number of experimental points into our training data, we can systematically correct biases inherited from purely theoretical calculations, thereby improving the transferability of the ML models.

Building on these established datasets, we aim to form collaborations with laboratories that specialize in advanced characterization—particularly institutions with robust X-ray diffraction (XRD) and NMR facilities—to obtain targeted measurements for materials relevant to nuclear and mining chemistry. For instance, precise XRD lattice constants can be used to refine structural predictions, while solid-state NMR chemical shifts can improve descriptions of local bonding and hydrogen bonding networks. In practice, these observables will be integrated into the ML-driven parameter-tuning process: the computational pipeline (whether it be adaptive DFT or the amons approach) will incorporate experimental discrepancies into its loss function or fit procedure, effectively “learning” how to align predicted properties with measured data. By closing this loop between theory and experiment, our models will rapidly identify and correct systematic errors, such as underestimated band gaps or mislocalized defect states, ultimately yielding more accurate simulations of complex materials under realistic conditions.


%%\section{Detailed Methodological Approach}
%
%To achieve the objectives outlined in this proposal, a multi-tiered methodological strategy will be employed. The first component focuses on the integration of adaptive ML techniques with dispersion-corrected DFT. Neural networks will be trained on extensive databases of high-level DFT calculations with an emphasis on accurately capturing the interplay of intermolecular interactions. By incorporating uncertainty quantification into these models, it will be possible to assess the reliability of predictions and identify regions in chemical space that warrant further high-level computations or would be better served by experimental results. Iterative refinement will be achieved through continuous feedback from new experimental data, ensuring that the adaptive algorithms remain accurate and relevant.
%
%The second component involves the development of a robust framework for solid-state amons. Complex crystal structures will be systematically decomposed into a set of fundamental fragments using advanced clustering and graph-based analysis techniques. Each fragment will be parameterized via dispersion-corrected DFT, resulting in a modular library that encapsulates the essential chemical and physical characteristics of the full material. This toolkit will be integrated into the FHIaims electronic structure package to enable efficient, high-throughput screening of defect-rich and layered materials, facilitating rapid property estimation and rational design.
%
%The third methodological component is the realization of a self-driving laboratory environment. This will be accomplished by interfacing computational tools with off the shelf automated synthesis platforms such as those being developed by the ChemSpeed company—incorporating robotic synthesis equipment and in-line characterization techniques. The integration will allow for real-time validation of computational predictions, establishing a dynamic feedback loop that continuously refines both the models and experimental protocols. Bayesian optimization techniques will be utilized to identify and fill gaps in the existing data, dramatically reducing the time and cost associated with discovering new materials.
%
%\section{Significance, Impact, and Differentiation}
%
%The anticipated impacts of this research program are multifaceted. First, by developing a suite of computational tools that combines the accuracy of dispersion-corrected DFT with the efficiency of adaptive ML and high-throughput experimentation, the program will set a new benchmark for computational materials design. The methods developed here will be broadly applicable, enabling rapid and accurate predictions across a range of materials classes.
%Second, the expansion of the solid-state amons framework represents a novel approach to decomposing complex materials into manageable, chemically meaningful subunits, and will provide a route to study solid form systems of scales that otherwise have not been possible. This modular strategy will facilitate targeted investigations into the role of defects and interlayer interactions in determining material properties, particularly in systems where conventional materials chemistry methods have encountered significant challenges.
%
%Importantly, this research program distinguishes itself from my previous work by explicitly venturing into new territories and addressing critical technological challenges. While my earlier research focused primarily on molecular crystal structure prediction, the present proposal extends these methodologies into areas with profound implications for next-generation electronic and optoelectronic devices, and combines them with the growing field of AI. My focused interest in semiconducting, 2D, and defect-rich materials reflects a deliberate diversification that positions the program at the forefront of computational materials chemistry. Moreover, the integration of computational methods with a self-driving laboratory marks a significant step toward bridging the gap between theoretical predictions and experimental realization, thereby accelerating the pace of materials discovery.
%


%% ------------
%% References
%% NSERC applications require the references to be in a distinct PDF,
%% uploaded separately from the main text. The hack I suggest is to
%% write the document *with* the references at the end (like a normal
%% article, and then to split it in two manually (for example, using pdftk).
%% For Linux systems, the script split-application.sh bundled with this
%% package does exactly that.
%% ------------
\newpage
\pagestyle{plain}
\bibliographystyle{unsrt}
\bibliography{example}

\end{document}
%% :folding=explicit:maxLineLen=76:wrap=soft:
